Findings of the Association for Computational Linguistics: EMNLP 2025, pages 15273–15291
November 4-9, 2025 ©2025 Association for Computational LinguisticsAdapting Large Language Models for Character-based Augmentative and
Alternative Communication
Dylan Gaines1,2* and Keith Vertanen1*
1Department of Computer Science, Michigan Technological University, Houghton, MI, USA
2Department of Computer Science, Kennesaw State University, Marietta, GA, USA
dgaine20@kennesaw.edu, vertanen@mtu.edu
Abstract
Users of Augmentative and Alternative Communication (AAC) may write letter-by-letter
via an interface that uses a character language
model. However, most state-of-the-art large
pretrained language models predict subword
tokens of variable length. We investigate how
to practically use such models to make accurate and efficient character predictions. Our
algorithm for producing character predictions
from a subword large language model (LLM)
provides more accurate predictions than using
a classification layer, a byte-level LLM, or an
n-gram model. Additionally, we investigate a
domain adaptation procedure based on a large
dataset of sentences we curated based on scoring how useful each sentence might be for spoken or written AAC communication. We find
our procedure further improves model performance on simple, conversational text.
1 Introduction
Augmentative and Alternative Communication
(AAC) devices allow non-speaking individuals to
communicate in face-to-face conversations and
asynchronously via messaging systems such as
email. Depending on their disability, AAC users
may be unable to use a conventional computer keyboard and may instead rely on alternative input
methods such as eye tracking or triggering a single switch (e.g. by twitching a muscle or puffing
air). Unfortunately, such input methods can make
writing slow and tiring. The writing speed of nonspeaking AAC users varies, but is frequently less
than 10 words per minute (WPM), with an average
closer to 5 WPM (Polacek et al., 2017).
While some AAC users may engage in long conversations, communications are often shorter and
more isolated in nature due to the input rate limitation. AAC users with severe motor impairments
may write letter-by-letter using an interface such
*These authors contributed equally to this work.
as Dasher (Ward and MacKay, 2002), the RSVP
Keyboard (Orhan et al., 2012), or Nomon (Broderick and MacKay, 2009). These AAC interfaces
all require the probability distribution over the next
character based on a user’s previous text.
We wanted to know if recent advances in pretrained large language models (LLMs) really help
with the “simple” case of predicting the next letter
in a sentence with only that sentence as context. Numerous papers have shown that ever-larger models
such as Llama 3 (Grattafiori et al., 2024), Mistral
7B (Jiang et al., 2023), and GPT-4 (OpenAI et al.,
2024) provide gains on more “complex” natural
language processing tasks such as question answering, sentiment analysis, and machine translation.
Many large transformer models do not operate
on a character level. Rather, they use a fixed vocabulary of subword tokens when making their predictions. These subword tokens can range anywhere
from a single character to a full word in length.
These tokens are learned during the training process and stem from frequently seen sequences of
characters. As we detail later, this subword tokenization makes it difficult to obtain single character predictions from many LLMs.
In this work we make four contributions: (1)
an algorithm that produces a probability distribution over single characters from a subword token
language model, (2) offline, novel test sets of AAClike text that can be used to evaluate LLMs without
fear of their inclusion in the models’ training data,
(3) a classification of existing large (87 B words)
training corpora based on their similarity to AAClike text, and (4) an empirical evaluation and comparison of different methodologies for obtaining
single-character predictions using LLMs.
We have released our algorithm’s source code,
our classified data sets, our classification model,
and our best language models as part of this work.1
1https://osf.io/ajm7t/
15273This paper is structured as follows: first, we motivate why character predictions are important for
AAC text entry (Section 2) and describe our algorithm for making such predictions using a subword
tokenized LLM (Section 3). Next, we find the best
starting model, testing performance on a private
collection of AAC-like text (Section 4). We then
explore adapting the LLM using small in-domain
sets and using text classified from much larger outof-domain sets (Section 5). Next, we compare our
search algorithm against other possible approaches
(Section 6). Finally, we show our results on five test
sets (Section 7) before wrapping up with discussion
(Section 8) and conclusions (Section 9).
2 Related work
2.1 Language modeling for AAC
Language models have often been studied in AAC
in hopes of improving communication speed. Work
has focused on four main topics: word completions
and predictions, sentence and phrase retrieval, abbreviation expansion, and domain adaptation.
Word completions and predictions. Word completions attempt to predict a user’s intended current
word before they finish typing it, while word predictions attempt to predict subsequent words based
on the likelihood of each word in the context of
the user’s sentence. Prior work by Kristensson and
Müllners (2021) showed that word predictions can
increase or decrease a user’s entry rate depending
on how and when predictions are displayed and the
user’s typing strategy. Yusufali et al. (2023, 2024)
used LLMs to generate word predictions for AAC
devices and showed entry rates up to 30.4 WPM.
Sentence and phrase retrieval. Retrieval techniques allow users to store sentences or phrases that
they frequently use (Todman et al., 2008). These
messages can then be retrieved by the user by entering specific keywords (Langer and Hickey, 1998)
or matching context tags (Kristensson et al., 2020).
The drawback to retrieval is utterances must be prestored by the user in their device, so they are not as
adaptable to different conversational situations.
Abbreviation expansion. Recent advances in
LLMs have allowed the development of more dynamic abbreviation expansion interfaces. Cai et al.
(2022, 2023) developed a system where users abbreviate each word to its first letter. The system
leverages the context of the utterance and an LLM
to decode the intended text. KWickChat (Shen
et al., 2022) allows the user to enter a few keywords and uses an LLM to generate expanded utterances based on the provided keywords, some
contextual details about the user, and the dialogue
from their conversational partner. Valencia et al.
(2023) also developed an abbreviation expansion
system that expands single words into complete
sentences using conversational context and background information about the user.
Domain adaptation. The type of text an AAC
user may want to write can differ from typical
sources of language model training data (e.g. text
from web pages). Vertanen and Kristensson (2011)
used a collection of crowdsourced fictional AAC
messages to filter sentences from large social media datasets using cross-entropy difference (Moore
and Lewis, 2010). Their optimized n-gram language models improved perplexity and potential
keystroke savings on AAC-like test sets. Adhikary
et al. (2021) used the same crowdsourced collection to select sentences from text scraped from the
web. They found using a BERT-based classifier
to select sentences slightly outperformed selection
via cross-entropy difference.
2.2 Applications for character predictions
While word and phrase-based language modeling
approaches can help accelerate input with many
AAC devices, some devices rely primarily on
character-at-a-time input. In these systems, it may
take several user interactions to input a single character due to a user’s noisy input signal. Character
predictions from a language model can help accelerate input of individual characters in such systems.
Brain-computer interfaces. P300 spellers (Farwell and Donchin, 1988) present visual stimuli to
users and measure their response using EEG electrodes. For example, the RSVP Keyboard (Orhan
et al., 2012) displays a rapid series of characters to
users. Using RSVP, it can take several series worth
of EEG responses to infer a user’s desired character.
RSVP uses a character n-gram language model in
combination with the EEG evidence to reduce the
number of series needed. We hypothesize an LLM
will provide more accurate probabilities and reduce
the number of series needed.
Nomon. Nomon (Broderick and MacKay, 2009;
Bonaker et al., 2022b,a) allows users to select
words or characters by activating a switch when the
rotating clock corresponding to their target reaches
noon. It can take several clicks to make a single
selection since the system needs to differentiate
between similar clocks. Since Nomon provides15274both word and character selection options, it needs
a language model that can evaluate the likelihood
of each possibility to help accelerate writing.
Dasher. Dasher (Ward and MacKay, 2002) inherently relies on character-level language model
predictions. In Dasher, the letters of the alphabet
are displayed vertically on the right side of the
interface. To enter text, users move their cursor
through the sequence of nested letter boxes corresponding to their desired text. Dasher can be
controlled by various devices such an eye-tracker,
mouse, or touchscreen. As users progress through
their text, more likely letter options are displayed
as larger targets, making it easier to select them.
Dasher uses the prediction by partial match
(PPM) algorithm (Cleary and Witten, 1984). PPM
is similar to an n-gram model but has a data structure making it easy to continually adapt to a user’s
writing. We hope to improve on this using LLMs.
However, a key challenge is generating predictions
quickly enough to support Dasher’s visualization
of the many possible future character sequences.
3 Character prediction algorithm
As we discussed previously, many modern large
transformer models use subword tokens instead of
single characters or full words. The first step in using a subword token language model is to convert
the context into a sequence of tokens via tokenization. Each model comes with its own tokenizer,
which uses a greedy algorithm, such as byte pair
encoding (Sennrich et al., 2016), to segment any
given text into tokens in its model’s vocabulary.
This subword tokenization complicates the process of obtaining a distribution of character probabilities since adding another character to the end
of a section of text could (1) begin a new token
containing only the new character, (2) extend the
existing most recent token with the new character,
or (3) rearrange the optimal tokenization of the
most recent word. For the last case, take for example writing the word “yesterday”. If the user has
written “yeste” thus far, the most likely tokenization may be into “y-este” (where - denotes token
boundaries). Adding an “r” could result in the most
likely tokenization changing to “yes-ter”.
To properly consider each of these possibilities,
we must remove one or more tokens from the end
of the context to allow the model the option of extending it, as opposed to only adding a new token.
Removing more tokens allows the model more flexibility to rearrange the optimal tokenization and in
our experience leads to better predictions, but at the
cost of a higher prediction time. In this work, we
remove all tokens in the current in-progress word to
allow maximal flexibility in tokenizing the partial
word. Going further back would provide no benefit
because, in all models we used, the space character
only appears at the start of tokens and thus tokens
cannot span word boundaries.
We then ask the model to estimate the likelihood of the next token and evaluate any token that
matches our context. For efficiency, we only track
a certain number of hypotheses at a time, known
as the beam width. We continue expanding each
hypothesis until it has regenerated the removed
context tokens plus at least one character. We then
store the likelihood for each final prediction in a
list based on the character that directly follows the
context. We continue this process until there are
no more hypotheses to extend, or, for efficiency,
the number of completed hypotheses exceeds a
max completed parameter. Then we sum the likelihoods stored for each character in our symbol set
and normalize to sum to one, arriving at a distribution over our symbol set. See Appendix C for our
pseudocode search algorithm or our supplementary
materials for a Python implementation.
To give a concrete example of the benefit of
removing tokens from the end of the context, we
further examine the “yeste” example. If we do
not remove any context, and simply marginalize
the probabilities over the first character in each
token, the probability of the next letter being R is
0.0047. Using our search algorithm with context
removal, this jumps to 0.9999. We suspect this is
due to the training process of the model; during
training it doesn’t frequently encounter fragmented
words and is allowed to tokenize the full word in the
most optimal way. By not removing any context,
we force the model to infer text with suboptimal
tokenizations, causing it to struggle.
4 Pretrained model performance
How do pretrained transformer language models
perform on AAC-like text? In this section, we
focus on LLMs using subword tokenization.
4.1 Private AAC-like test data
Unfortunately, actual text from AAC users is difficult to source due to practical, ethical, and privacy
reasons. There are a number of datasets that have15275been used to approximate AAC communications
such as voice banking messages (Costello, 2014),
Switchboard (Godfrey et al., 1992), DailyDialog
(Li et al., 2017), COMM2 (Vertanen, 2013), and
Turk dialogues (Vertanen, 2017). However, these
datasets are available online and may have been
used in training the LLMs we want to evaluate.
As we wish to fairly compare performance between different LLMs, as well as between LLMs
and an n-gram language model optimized for AAClike text, we created two private tests sets, one for
informal written communications (i.e. similar to
mobile messaging), and one for person-to-person
communications. We denoted the two distinct genres as written and spoken respectively.2 We sourced
our text from compositions made by workers on
Amazon Mechanical Turk in response to various
prompts. For details of how we constructed our test
sets, including example sentences, see Appendix A.
We split each genre in half, creating a written dev
set (8.6 K words), a written test set (9.0 K words),
a spoken dev set (11.3 K words), and a spoken test
set (10.8 K words). We use the dev sets for most of
our experiments, reserving the test sets for a final
evaluation (Section 7). Our final evaluation also
uses text written by an AAC user, messages written
by people with ALS for voice banking (Costello,
2014), and responses to common conversational
situations (Vertanen, 2013).
4.2 Perplexity experiments
We converted our private test sets to lowercase and
removed anything not an English letter, apostrophe,
or space. We calculated per-character perplexity
(including spaces) on our dev sets without any sentence end token and using the sentence start token
appropriate for the model under test.
We used 16-bit floating point inference for the
LLMs. We used our search algorithm (Section 3)
for LLM predictions. We found a search beam of 8
and a max completed stopping criteria of 32 K provided a perplexity close to the perplexity without
pruning while significantly speeding inference.
We compared against a baseline of a 12-gram
language model optimized for AAC-like text, denoted AAC n-gram.3 This model was trained using
Witten-Bell smoothing on 21 B characters from
Common Crawl, Reddit, movie subtitles, and Twitter. Training sentences were selected via cross-
2We intend to provide access to researchers on the condition they not leak the text online or to hosted models.
3https://osf.io/c6mnz
entropy difference selection (Moore and Lewis,
2010) using in-domain models trained on DailyDialog (Li et al., 2017), short emails, and AAC-like
messages (Vertanen and Kristensson, 2011).
As shown in Table 1, per-character perplexity
decreased with increasing model size, though gains
plateaued for the OPT model family (Zhang et al.,
2022) at around 1.3 B parameters. Thus, even given
the “easy” nature of our task (i.e. predicting the
next character within an isolated sentence), bigger
models were often better. While the n-gram model
optimized for AAC did better than DistilGPT with
a similar number of parameters, larger LLMs eventually outperformed it despite the LLMs not being
fine-tuned on AAC-like text.
The OPT models outperformed the GPT-2 models (Radford et al., 2018) at similar model sizes. We
also tested the recent Llama 3 model (Grattafiori
et al., 2024) that was trained on substantially more
data than the other models (15 T tokens), but found
it did no better than the OPT models.
Determining the cause of the perplexity differences between model families is difficult due to the
many differences in the underlying LLM architectures, training data, and training procedures. But
given that the performance of n-gram models on
AAC-like text have been shown to depend on the
training data (Vertanen and Kristensson, 2011), we
conjecture the difference may be due to the amount
of well-matched English text versus other types of
text (e.g. other natural or programming languages).
Our n-gram model inferences using KenLM
(Heafield, 2011) were fast at 0.1 ms per prediction
(using a CPU). Despite substantial performance
optimization efforts as well as using an NVIDIA
A100 GPU for inference, our search algorithm was
much slower at 55–120 ms per prediction. In general, larger models took longer. We will investigate
other methods for faster inference in Section 6.
5 Optimizing language models for AAC
The LLMs in the previous section were all trained
on general text from the web. As with previous
work on n-gram language models (Vertanen and
Kristensson, 2011; Adhikary et al., 2021), we may
be able to improve predictions on AAC-like text
by focusing the model on this style of text. We
selected the opt-350m model for further refinement
as it provided the majority of perplexity gains while
being small enough to plausibly be used on an enduser’s device. For privacy and latency reasons,15276Params (M) Perplexity Time (ms)
AAC n-gram 70 2.55 0.1
distilgpt2 82 2.71 55.0
gpt2 124 2.54 60.0
gpt2-medium 355 2.45 70.1
gpt2-xl 1558 2.38 91.7
opt-125m 125 2.39 64.3
opt-350m 331 2.31 78.6
opt-1.3b 1316 2.22 77.9
opt-2.7b 2652 2.21 86.8
opt-6.7b 6658 2.20 89.5
opt-13b 12853 2.23 108.8
Meta-Llama-3.1-8B 8030 2.31 119.9
Meta-Llama-3.2-1B 1236 2.37 97.0
Meta-Llama-3.2-3B 3213 2.33 115.0
Table 1: Model parameters, per-character perplexity,
and prediction time of pretrained LLMs on our dev sets
(average of the written and spoken sets).
users may prefer to use a local model.
Our results in Table 1 were on lowercase versions of our test sentences. Our search algorithm
sums over all matching following tokens ignoring
case. However, we conditioned predictions on all
lowercase context. This simulates an AAC interface where the user does not bother to denote case,
or when casing is added later in some way. We
found we could slightly improve perplexity on our
dev sets using opt-350m from 2.31 to 2.30 by upper casing the first letter of the sentence and the
words: I, I’m, I’ll, I’ve, and I’d. We found this
simple approach performed similar to using the
human-supplied case. We use this automatic casing method for the remainder of our experiments.
5.1 In-domain datasets
To approximate written AAC, we used sentences
from forum posts made by people using a mobile
device (Vertanen and Kristensson, 2021). We denote this dataset as Forum. Our rationale for using
this dataset was that it contains a rich set of topics
while also having a bias towards concise writing.
To approximate spoken AAC, we used the DailyDialog corpus (Li et al., 2017) and the BOLT
SMS/Chat corpus (Song et al., 2014). We denote
these datasets as Daily and BOLT respectively. Our
rationale for choosing these datasets was that while
the text was actually written and not spoken, it focuses on person-to-person dialogues. Compared
to other datasets based on spoken dialogues such
as Switchboard (Godfrey et al., 1992), the text is
relatively clean and does not contain spoken speech
artifacts. Daily had 46 K sentences (356 K words),
BOLT had 271 K sentences (1.9 M words), and Forum had 888 K sentences (10.1 M words).
5.2 Large datasets
We selected two large corpora to mine data from:
the Colossal Clean Crawled Corpus (Raffel et al.,
2020) (ODC-By License) and the OpenSubtitles2016 movie subtitle corpus (Lison and Tiedemann, 2016) (License unknown). We denote these
datasets as C4 and Subtitle respectively. This selection was guided by results4 showing that out
of 14 possible training sources, these two corpora performed the best on modeling AAC-like
text. Further, we think C4 provides text that wellapproximates a written style of communication
while Subtitle represents a more spoken style.
We filtered C4 and Subtitle to sentences that
started with a capital letter, had both uppercase and
lowercase characters, and ended with sentence end
punctuation. We dropped sentences where more
than 20% of the words were out-of-vocabulary
with respect to a list of 860 K words obtained from
human-edited dictionaries. Throughout this work,
we trained and tested on isolated sentences. We
removed repeated identical sentences in Subtitle.
We reserved 5% of the data as a dev set and 5% as
a test set. After filtering, our C4 training set had
87.5 B words and Subtitle had 528 M words.
5.3 Sentence classification
Similar to past work (Vertanen and Kristensson,
2011; Adhikary et al., 2021), we wanted to focus
our models on sentences similar in style to the text
AAC users may want to write. We did this by training a three-way classifier on top of DeBERTaV3
(He et al., 2021). For the out-of-domain class, we
used a newswire corpus (Chelba et al., 2013) which
we denote as News. For the written in-domain class,
we used Forum. For the spoken in-domain class,
we used a combination of Daily and BOLT.
Daily was the smallest with 46 K sentences. We
used NLPAug (Ma, 2019) to add 232 K sentences,
changing one or two words (selected at random) in
the original sentences. Words were changed using
a distilled version of RoBERTa (Liu et al., 2019).
This brought Daily to the same size (279 K sentences) as BOLT. We took the first 558 K sentences
from Forum and News to create an equal amount of
data for each of our three classes (in-domain written, in-domain spoken, and out-of-domain). Dur-
4https://osf.io/ajm7t/wiki/lm_char_dec19/15277ing classification, we lowercased sentences and
removed characters aside from A–Z, apostrophe,
and space. We optimized hyperparameters as described in Appendix B.1. We used early stopping
during training. Our final model had a classification
accuracy of 85.7% on unseen sentences.
5.4 Domain adaptation procedure
We scored each sentence in C4 and Subtitle using
our classifier. To adapt the opt-350m LLM, we
selected sentences with a written or spoken probability that met a threshold. To efficiently determine
which threshold to use, we trained 12-gram models
from scratch across a broad range of thresholds
and found we obtained the majority of perplexity
gains with a threshold of 0.90 for C4 and 0.75 for
Subtitle. Using these thresholds, our adaptation
datasets consisted of 4.0 B words of C4 and 130 M
words of Subtitle. With these thresholds as a starting point, we then trained opt-350m models with a
narrower range of thresholds centered on 0.90 for
C4 and 0.75 for Subtitle. The results, detailed in
Appendix B.3, confirmed the thresholds we chose.
We first tested fine-tuning all the model’s parameters on a combined adaptation set consisting of C4
(threshold 0.90), Subtitle (threshold 0.75), BOLT,
and Daily. As shown in Table 2, this reduced the
perplexity on our dev sets from 2.30 to 2.19. Next,
we explored a curriculum adaptation process consisting of a sequence of separate fine-tuning runs.
Each run was on one of our five datasets: C4, Subtitle, Daily, BOLT, or Forum. Before each finetuning, we optimized hyperparameters using the
model from the previous step.
Given the sizes of our five datasets, we opted to
first always fine-tune on C4 followed by Subtitle.
We tested all possible orders of further fine-tuning
on Daily, BOLT, and Forum. Fine-tuning on Forum
did not help in any order, and thus we excluded it
from the curriculum. We found fine-tuning on C4,
Subtitle, BOLT, and then Daily was the best order,
improving perplexity to 2.10. Combining BOLT
and Daily in a single final step did slightly better
at 2.09. However, simply fine-tuning using the
combined BOLT and Daily data did slightly better
still at 2.07. While the classified C4 and Subtitle
datasets did not prove useful for improving the
subword LLM, they will be useful for creating or
improving other models (see Section 6).
Additionally, we fine-tuned opt-350m using
Low-Rank Adaptation (LoRA) (Hu et al., 2021),
which drastically reduces the number of trainable
Adaptation data Fine-tuning Perplexity
type steps Written Spoken Avg
— — — 2.41 2.19 2.30
C4+Sub+BOLT+Daily full 1 2.30 2.08 2.19
C4 )Sub )BOLT )Daily full 4 2.16 2.03 2.10
C4 )Sub )BOLT+Daily full 3 2.14 2.03 2.09
BOLT+Daily full 1 2.11 2.04 2.07
C4 )Sub )BOLT+Daily LoRA 3 2.17 2.06 2.11
BOLT+Daily LoRA 1 2.13 2.05 2.09
Table 2: Perplexity on the two dev sets using various
adaptation approaches on opt-350m. Subtitle (denoted
Sub) used a 0.75 threshold and C4 used a 0.90 threshold.
) denotes separate fine-tuning steps. + denotes datasets
merged in a single fine-tuning step. Best result in bold.
parameters for downstream tasks. As shown in
Table 2 (bottom), LoRA yields slightly higher perplexities compared to full parameter fine-tuning.
As such, we continued to use full fine-tuning for
the remainder of our experiments.
For additional results and hyperparameter values,
see Appendix B.2. We also explored other ways to
try and improve adaptation. For details of things
that did not work, see Appendix B.3.
6 Other prediction approaches
With opt-350m, our search algorithm (Section 3)
takes on average 79 ms to predict the next character
on an A100 GPU. While this may be fast enough
for many purposes, interfaces such as Dasher (Ward
and MacKay, 2002) or less capable devices may
require more efficient prediction. We explored several other methods for obtaining faster predictions.
6.1 N-gram model
Inference with an n-gram model was very fast at
0.1 ms (Table 1). To serve as a baseline, and to
provide an open training data n-gram model, we
created a character 12-gram n-gram mixture model
from C4 (threshold 0.90) and Subtitle (threshold
0.75) using SRILM (Stolcke, 2002). We created
separate C4 and Subtitle language models using
Witten-Bell smoothing and no count cutoffs. We
built a mixture model using linear interpolation
with weights optimized to minimize perplexity on
dev sentences from Forum, BOLT, and Daily. The
mixture weights were 0.84 (C4) and 0.16 (Subtitle).
The unpruned mixture model (denoted C4+Sub
n-gram full) was large with 5 B parameters and a
compressed size of 30 GB. We used entropy pruning (Stolcke, 1998) to to reduce the model to 71 M15278parameters and a compressed size of 488 MB (denoted C4+Sub n-gram). This was similar in size to
the AAC n-gram we previously compared with (Table 1). On our dev sets, our full and pruned models
had perplexities of 2.53 and 2.60 respectively. Our
pruned model’s perplexity was somewhat higher
than the 2.55 of the existing AAC n-gram model.
6.2 Byte model
The multilingual encoder-decoder ByT5 model
(Xue et al., 2022) eschewed the use of subword
tokenization and instead used a token set based on
individual characters. ByGPT5 (Belouadi and Eger,
2023) is a decoder-only version of ByT5 that can
make character predictions directly.
The average perplexity on our dev sets of the
base ByGPT5 model (289 M parameters) was 2.85.
We compared our two best adaptation methods we
previously found on opt-350m (Section 5.4) using
ByGPT5. First, we used a three step curriculum
approach: adapting on C4, then on Subtitle, and
finally on BOLT combined with Daily. This model
had a perplexity of 2.14.
Second, we fine-tuned the base model in a singlestep on BOLT combined with Daily. This results
in a higher perplexity of 2.27. We suspect due
to the multilingual nature of the base model, first
fine-tuning on C4 and Subtitle helped move the
model towards English. For results and hyperparameter values at each phase of the fine-tuning, see
Appendix B.4.
6.3 Classification
We also explored making character predictions by
adding a classification layer to the LLM. Instead of
using the model to generate the ensuing text from
some context, we asked it to produce a distribution
over labels – each label represented a character. We
built the model on top of the best domain-adapted
opt-350m model from Section 5.4. This retained
the model’s domain knowledge and focused additional training on the classification task.
To create training examples, we selected a random location within each sentence. We truncated
the sentence at that point and used the character
that would come next as the label. If the next character was not in our symbol set (A–Z, a–z, space,
and apostrophe), we chose another random location. We utilized the same three step curriculum
approach as the byte model; we first trained on C4,
then Subtitle, and finally on BOLT combined with
Daily. For the smaller BOLT and Daily sets, we
added a hyperparameter for the number training
examples generated per sentence. During training,
we updated the classification layer as well as all
other model parameters.
The classifier model did not perform as well as
we had hoped, achieving a perplexity of 2.47. By
comparison, training with BOLT and then Daily in
separate steps yielded a slightly higher perplexity
of 2.54. Training on just BOLT and Daily performed worse at 2.60. See Appendix B.5 for additional results and hyperparameter values.
7 Final evaluation
At this point, we have extensively used our private
dev sets to guide our process. We conducted a final evaluation with our unseen private test sets and
three additional test sets. Firstly, we created a test
set based on the text of an AAC user. This set consisted of 3.6 K sentences (65 K words) written by a
user of Dasher (Ward and MacKay, 2002). The text
was from public and private presentations given by
the user. We also added two public test sets: voice
banked messages of people with ALS (Costello,
2014), and the COMM2 set of responses to conversational situations (Vertanen, 2013).
As shown in Table 3 (top), our C4+Sub n-gram
model performed slightly worse with an average
test set perplexity of 2.58 compared to the AAC
n-gram model which averaged 2.54. Our unpruned
full model did slightly better at 2.53. All n-gram
models were outperformed by the LLMs, both with
and without domain adaptation.
As shown in Table 3 (middle), using the search
algorithm with opt-350m adapted in a single step
on the combination of BOLT and Daily performed
the best on two of the five test sets and the overall
average. It tied with the four-step curriculum on the
COMM2 set (2.01). The base opt-350m yielded the
lowest perplexity on the test set from the AAC user
(2.24), though the curriculum adapted opt-350m
models were not far behind (2.28–2.30).
Averaging the perplexity across all models, there
was a noticeable increase in perplexity between
our written (2.35) and spoken (2.24) test sets and
the presentations given by the AAC user (2.59).
The public test sets also produced lower average perplexities (2.28 for ALS bank and 2.21 for
COMM2). While we hoped our data selection approach using diverse text from mobile forum messages would allow our adapted models to perform
well on both simple everyday communications and15279Language Domain Character Written Spoken AAC user ALS bank COMM2 Average Time
model adaptation prediction (ppl) (ppl) (ppl) (ppl) (ppl) (ppl) (ms)
AAC n-gram - KenLM 2.56 2.51 2.89 2.33 2.41 2.54 0.1
C4+Sub n-gram - KenLM 2.62 2.54 2.86 2.43 2.45 2.58 0.1
C4+Sub n-gram full - KenLM 2.57 2.47 2.82 2.38 2.40 2.53 0.1
opt-350m - search alg. 2.40 2.15 2.24 2.39 2.20 2.28 79.1
opt-350m C4 )Sub )BOLT )Daily search alg. 2.14 2.00 2.29 2.14 2.01 2.12 78.9
opt-350m C4 )Sub )BOLT+Daily search alg. 2.12 2.00 2.30 2.14 2.02 2.12 78.2
opt-350m BOLT+Daily search alg. 2.10 2.01 2.28 2.13 2.01 2.11 78.0
ByGPT5 C4 )Sub )BOLT+Daily directly 2.20 2.04 2.57 2.18 2.06 2.21 11.6
opt-350m C4 )Sub )BOLT+Daily classifier 2.42 2.42 3.03 2.38 2.36 2.52 14.6
Table 3: Final evaluation of our various language models and methods for predicting the next character. ) denotes
separate fine-tuning steps. + denotes datasets merged in a single fine-tuning step. Inference times are the average of
per-character inferences over the test sets. The best result in each column is shown in bold.
Language model Domain adaptation Written Spoken AAC user ALS bank COMM2 Average
AAC n-gram - 59.7% 60.0% 55.3% 62.2% 61.0% 59.7%
C4+Sub n-gram - 59.2% 59.9% 56.4% 61.1% 61.2% 59.6%
C4+Sub n-gram full - 60.1% 61.4% 57.4% 61.8% 62.0% 60.5%
opt-350m - 61.7% 65.1% 63.1% 61.9% 64.2% 63.2%
opt-350m C4 )Sub )BOLT )Daily 64.6% 67.1% 62.5% 64.5% 66.4% 65.0%
opt-350m C4 )Sub )BOLT+Daily 64.9% 67.0% 62.3% 64.4% 66.3% 65.0%
opt-350m BOLT+Daily 65.2% 67.2% 62.7% 64.7% 66.3% 65.2%
ByGPT5 C4 )Sub )BOLT+Daily 64.3% 66.8% 61.6% 63.7% 65.9% 64.5%
Table 4: Keystroke savings for language models using an onscreen predictive keyboard with five word predictions.
The best result in each column is shown in bold.
more complex planned written communications,
at least for this one user, this did not work. As a
measure of text complexity, we calculated the number of characters per sentence. The AAC user’s
sentences were on average 94.4 characters per sentence compared to the much shorter sentences in
the written (31.4) and spoken (39.9) test sets.
Of the LLM approaches, the byte LLM provided
the fastest predictions at around 12 ms. This was
6.8 times faster than our search algorithm and 1.3
times faster than using a classification layer.
To demonstrate the potential impact in a realworld AAC interface, we simulated using each language model in a onscreen keyboard with five word
predictions. We simulated the keystrokes a hypothetical perfect user would save compared to typing
every character. Keystroke savings is calculated as:
ka − kp
ka
× 100%,
where ka is the number of keystrokes required
without word predictions and kp is the number
of keystrokes required with predictions. Higher
keystroke savings is better.
The keyboard was allowed to make up to five
word predictions. Predictions were made before
any characters were typed for the next target word
as well as during each subsequent prefix of the
target word. The simulated user made no typing errors. We did not filter the offered word predictions
based on whether the user did not choose a given
prediction for a shorter prefix for the current target.
If the target word appeared in the predictions, we
assumed one key press selected the target word and
added any following space.
As shown in Table 4, keystroke savings improvements followed the perplexity gains (Table 3).
Compared to the AAC n-gram, across the test sets
our best LLM provided a 5.5% absolute increase in
keystroke savings. To contextualize this, imagine
someone is using an eye-tracker and an onscreen
keyboard with a dwell time of one second. Assuming they make no errors and incur no other overheads, writing the average length sentence in our
spoken test set would take about 40 seconds without word predictions. Making optimal use of predictions provided by the AAC n-gram would take
about 16.0 seconds. Using our best model (opt350m adapted on BOLT+Daily), the user would
take about 13.1 seconds, 18% faster than using the
n-gram model for predictions.152808 Discussion
Most of our trained models had higher perplexities
on the AAC test set than on any of the other test sets.
The only exception was the base opt-350m model,
which performed worst on the written test set. We
believe this stems from the nature of the test sets.
The written and spoken sets were created by crowdsourced workers writing about topics such as the
weather or travel. As described in Appendix A.4,
workers were instructed to invent questions or statements they might include in a conversation with
another person. The ALS bank and COMM2 sets
contain text of a similar nature. This aligns closely
with the in-domain conversational datasets we used
in our adaptation process.
The AAC test set was derived from prepared presentations by an AAC user with ALS discussing
their life experiences. It covers more complex topics and uses more advanced language. Presentations are also not conversational in nature, which
is why we believe our curriculum adaptation process was not effective at improving the model’s
performance on the AAC test set. We also note that
AAC users are diverse in nature and our domain
adaptation may be advantageous to other users or
even to this user in other communication situations.
It may also be advantageous to consider adapting
the language model more directly to a particular
user in lieu of, or in addition to, domain adaptation.
We obtained the lowest perplexities using a subword tokenized LLM and our search algorithm.
While we investigated a multi-step curriculum to
adapt the opt-350m model using sentences classified from large out-of-domain datasets, a simple
fine-tuning using in-domain datasets worked better.
The main drawback of the search algorithm is its
inference time. If faster predictions are required for
a particular interface or device, a byte or classification model may be preferred. However, currently
subword LLMs are much more popular than byte
LLMs. Further, the search approach can be used
directly on an LLM without task-specific training.
Adding a classification layer to the domainadapted opt-350m model performed worse than
any of the other LLM-based approaches, including
the base opt-350m model before domain adaptation.
We have a few theories as to why this may have
been the case. Our classification training examples
chose a random location in each sentence, so many
of the sentences included in these further training
steps were likely fragmented, which may have affected earlier layers in the model. To investigate
this theory, we could freeze the lower model layers
during the classification training. This might help
to preserve the domain adaptation and prior model
while also tuning the new classification layer.
It is also possible that the subword nature of the
model prevented it from fully learning the character classification task. As discussed in Section 3,
adding a new character can rearrange a word’s optimal tokenization. Because the model was originally
trained to predict the next token in a sequence, it
may not have considered all the different possible
tokenizations as our search algorithm does.
9 Conclusions
We investigated how to leverage large transformer
language models to aid non-speaking individuals
using Augmentative and Alternative Communication (AAC) devices. In particular, we focused on
how to generate the character predictions needed
by interfaces where users write one letter-at-a-time.
We presented an algorithm for obtaining character
probabilities from LLMs using subword tokenization, and compared our algorithm to using a bytelevel LLM or adding a classification layer on top of
a subword LLM. We detail how we classified sentences in large corpora of web and movies subtitles
text based on how similar sentences were to written and spoken communications. We investigated
a multi-step fine-tuning process using our scored
sentences as well as our smaller in-domain corpora
to adapt LLMs to our target domain.
We found that we were able to produce the lowest per-character perplexities using our search algorithm with the subword opt-350m model. We
found that our domain adaptation curriculum was
effective at improving the model performance for
simple, conversational text, but did not improve
performance for more complex prepared presentations from a single AAC user. Though it offered
the best perplexity, our subword to character algorithm took almost seven times longer to produce
character predictions than a native byte-level LLM.
To aid future research, we have made our unique
datasets and models available.5 This includes the
model for classifying sentences as written or spoken, scored C4 and Subtitle sentences, n-gram models trained on the classified sentences, fine-tuned
opt-350m model, fine-tuned ByGPT5 model, and
letter classification model.
5https://osf.io/ajm7t/15281Limitations
The ultimate goal of this work was to adapt large
language models for use in AAC devices. With that
goal in mind, it is a limitation that we did not have
training data from actual AAC users. As we discussed in Section 4.1, such data is difficult to come
by, so the datasets we created and used in this work
aimed to approximate AAC-like text. While we
think our models can provide an improved initial
experience for AAC text entry interfaces, we think
much more substantial improvements are likely
possible if models are adapted using the actual
text of a given user. Doing this in a practical and
privacy-preserving way will need further research.
The experiments we conducted as part of this
research consisted solely of the offline evaluation
of our models with fixed datasets. While the perplexity values and keystroke savings we report are
generally considered to be indicative of a model’s
performance, they may not show the full extent of
the impact on an actual user entering text. Text
input is a complex task, especially when performed
through an AAC device. Further research and user
evaluation is required to fully understand the benefit or detriment of our tuned language models on
the text entry performance of AAC users.
Some recent work has explored freezing model
layers that have relatively low training loss in order
to focus fine-tuning on improving underperforming layers (Yusufali et al., 2024). This was not a
technique that we were able to explore in the scope
of this work. However, we may have been able
to improve the performance of the classification
model by freezing the lower model layers that were
already domain adapted and focusing the additional
training on the classification layer specifically.
We only tested our search algorithm with a few
model architectures. Of particular concern is our
handling of token removal. Because the tokenizer
used by the models in this work includes a space at
the beginning of many of its tokens, our algorithm
removes the end of the context up to and including
the last space. Other models may have tokenizers
that handle space differently. Our algorithm may
need adjustment to work with such models.
The main advantage of our search algorithm
compared to the byte model is that it works out
of the box with the most common LLMs which use
subword tokenization. However, the main disadvantage is obtaining the distribution over the next
character takes substantially longer. The search
itself in many cases needs to make several inferences on the GPU so we can expect it to take several times as long as a byte LLM or classification
head-based approach which can make the character
prediction in a single inference. During the development of our search algorithm, we conducted
multiple rounds of performance optimization. We
introduced various changes to the algorithm that, in
the end, resulted in performance that was several orders of magnitude faster than our initial prototype.
However, it is possible additional efforts could further improve performance.
Currently, our search algorithm is optimized for
inference on a GPU and not a CPU. This could
limit the AAC devices our algorithm can practically be deployed on. However, for many AAC
interfaces, such as those based on switch or brain
input, waiting a tenth of a second for a prediction
is likely acceptable and would not unduly impact
user experience.
We primarily focused on fine-tuning approaches
to domain adaptation and the task of classifying
the next character. In-context learning (ICL) has
emerged in recent years as an alternative for finetuning (Dong et al., 2024). With in-context learning, a few examples of the task or domain are provided to the language model along with the query.
The model takes the examples into account when
producing its response, instead of requiring a separate training procedure. While ICL has shown
promise, such emergent abilities may require large
model sizes (Wei et al., 2022). For privacy reasons,
our focus here was on modest sized models that
could be plausibly used and fine-tuned on an end
user’s device. We were not able to explore the feasibility of ICL using modest sized models in the
scope of this paper, and leave this as future work.
Ethical considerations
Our fine-tuning process presented in this work relied heavily upon the Forum (Vertanen and Kristensson, 2021), BOLT (Song et al., 2014), and DailyDialog (Li et al., 2017) datasets. While we did
introduce the additional larger datasets C4 (Raffel
et al., 2020) and Subtitle (Lison and Tiedemann,
2016), we filtered those sets to sentences that were
similar to our in-domain sets. This achieved our intended result of improving the performance of our
character prediction models. However, any biases
present in Forum, BOLT, and Daily may also be
reflected in the text selected from C4 and Subtitle.15282This could lead to predictions made by text entry
interfaces relying on our models to skew towards
those biases.
When developing AAC interfaces that leverage
language models to accelerate user input, it is crucial to support user autonomy, i.e. allowing users
to express exactly what they want and not just what
is probable under a language model. We developed
our models to assist users in entering text based
on the text that is most likely. However, there may
be cases where users wish to enter text that is not
likely, or that the model has not seen before. We
encourage anyone who uses our models to design
their interfaces such that users can control to what
degree, if any, the language model contributes to
the produced text.
Acknowledgments
We thank the AAC user who provided their text
for use in our final evaluation. We thank Soufia
Bahmani for her work on the subword LLM word
prediction algorithm used in Table 4. This work
was funded by the National Institutes of Health
/ National Institute on Deafness and Other Communication Disorders (R01DC009834) and by the
National Science Foundation (IIS-1750193 and IIS2402876). Any opinions presented in this work are
those of the authors and do not reflect the opinions
of our funding agencies.
References
Jiban Adhikary, Jamie Berger, and Keith Vertanen. 2021.
Accelerating Text Communication via Abbreviated
Sentence Input. In Proceedings of the Joint Conference of the 59th Annual Meeting of the Association
for Computational Linguistics and the 11th International Joint Conference on Natural Language Processing, pages 6574–6588.
Takuya Akiba, Shotaro Sano, Toshihiko Yanase, Takeru
Ohta, and Masanori Koyama. 2019. Optuna: A nextgeneration hyperparameter optimization framework.
In Proceedings of the 25th ACM SIGKDD International Conference on Knowledge Discovery and Data
Mining.
Jonas Belouadi and Steffen Eger. 2023. Bygpt5: End-toend style-conditioned poetry generation with tokenfree language models. In Proceedings of the 61st
Annual Meeting of the Association for Computational Linguistics (Volume 1: Long Papers), page
7364–7381. Association for Computational Linguistics.
Nicholas Bonaker, Emli-Mari Nel, Keith Vertanen, and
Tamara Broderick. 2022a. Demonstrating nomon: A
flexible interface for noisy single-switch users. In
Extended Abstracts of the 2022 CHI Conference on
Human Factors in Computing Systems, CHI EA ’22,
page 1–4, New York, NY, USA. Association for Computing Machinery.
Nicholas Ryan Bonaker, Emli-Mari Nel, Keith Vertanen, and Tamara Broderick. 2022b. A performance
evaluation of nomon: A flexible interface for noisy
single-switch users. In CHI Conference on Human
Factors in Computing Systems, CHI ’22, New York,
NY, USA. Association for Computing Machinery.
Tamara Broderick and David J. C. MacKay. 2009. Fast
and flexible selection with a single switch. PLoS
ONE, 4(10):e7481.
Shanqing Cai, Subhashini Venugopalan, Katie Seaver,
Xiang Xiao, Katrin Tomanek, Sri Jalasutram, Meredith Ringel Morris, Shaun Kane, Ajit Narayanan,
Robert L. MacDonald, Emily Kornman, Daniel
Vance, Blair Casey, Steve M. Gleason, Philip Q. Nelson, and Michael P. Brenner. 2023. Using large
language models to accelerate communication for
users with severe motor impairments. axXiv preprint
arXiv:2312.01532.
Shanqing Cai, Subhashini Venugopalan, Katrin
Tomanek, Ajit Narayanan, Meredith Morris, and
Michael Brenner. 2022. Context-aware abbreviation expansion using large language models. In Proceedings of the 2022 Conference of the North American Chapter of the Association for Computational
Linguistics: Human Language Technologies, page
1261–1275, Seattle, United States. Association for
Computational Linguistics.
Ciprian Chelba, Tomáš Mikolov, Mike Schuster, Qi Ge,
Thorsten Brants, Phillipp Koehn, and Tony Robinson. 2013. One billion word benchmark for measuring progress in statistical language modeling. arXiv
preprint arXiv:1312.3005.
J. Cleary and I. Witten. 1984. Data compression using
adaptive coding and partial string matching. IEEE
Transactions on Communications, 32(4):396–402.
John M Costello. 2014. Message banking, voice banking and legacy messages. Boston Children’s Hospital,
Boston, MA.
Qingxiu Dong, Lei Li, Damai Dai, Ce Zheng, Jingyuan
Ma, Rui Li, Heming Xia, Jingjing Xu, Zhiyong Wu,
Tianyu Liu, Baobao Chang, Xu Sun, Lei Li, and
Zhifang Sui. 2024. A survey on in-context learning.
L.A. Farwell and E. Donchin. 1988. Talking off the top
of your head: toward a mental prosthesis utilizing
event-related brain potentials. Electroencephalography and clinical neurophysiology, 70(6):510–523.
J.J. Godfrey, E.C. Holliman, and J. McDaniel. 1992.
SWITCHBOARD: Telephone speech corpus for research and development. Proceedings of the IEEE
Conference on Acoustics, Speech, and Signal Processing, pages 517–520.15283Aaron Grattafiori, Abhimanyu Dubey, Abhinav Jauhri,
and Abhinav Pandey et al. 2024. The Llama 3 herd
of models.
Pengcheng He, Jianfeng Gao, and Weizhu Chen. 2021.
Debertav3: Improving deberta using electra-style pretraining with gradient-disentangled embedding sharing.
Kenneth Heafield. 2011. KenLM: faster and smaller language model queries. In Proceedings of the EMNLP
2011 Sixth Workshop on Statistical Machine Translation, pages 187–197, Edinburgh, Scotland, United
Kingdom.
Edward J. Hu, Yelong Shen, Phillip Wallis, Zeyuan
Allen-Zhu, Yuanzhi Li, Shean Wang, Lu Wang,
and Weizhu Chen. 2021. Lora: Low-rank adaptation of large language models. arXiv preprint
arXiv:2106.09685. ArXiv:2106.09685 [cs].
Shengding Hu, Yuge Tu, Xu Han, Chaoqun He, Ganqu
Cui, Xiang Long, Zhi Zheng, Yewei Fang, Yuxiang Huang, Weilin Zhao, et al. 2024. Minicpm:
Unveiling the potential of small language models
with scalable training strategies. arXiv preprint
arXiv:2404.06395.
Albert Q. Jiang, Alexandre Sablayrolles, Arthur Mensch, Chris Bamford, Devendra Singh Chaplot, Diego
de las Casas, Florian Bressand, Gianna Lengyel, Guillaume Lample, Lucile Saulnier, Lélio Renard Lavaud,
Marie-Anne Lachaux, Pierre Stock, Teven Le Scao,
Thibaut Lavril, Thomas Wang, Timothée Lacroix,
and William El Sayed. 2023. Mistral 7b. arXiv
preprint arXiv:2310.06825. ArXiv:2310.06825 [cs].
Per Ola Kristensson, James Lilley, Rolf Black, and Annalu Waller. 2020. A design engineering approach
for quantitatively exploring context-aware sentence
retrieval for nonspeaking individuals with motor disabilities. In Proceedings of the 2020 CHI Conference
on Human Factors in Computing Systems, page 1–11,
Honolulu HI USA. ACM.
Per Ola Kristensson and Thomas Müllners. 2021. Design and analysis of intelligent text entry systems
with function structure models and envelope analysis. In Proceedings of the 2021 CHI Conference on
Human Factors in Computing Systems, page 1–12,
Yokohama Japan. ACM.
Stefan Langer and Marianne Hickey. 1998. Using semantic lexicons for full text message retrieval in a
communication aid. Natural Language Engineering,
4(1):41–55.
Yanran Li, Hui Su, Xiaoyu Shen, Wenjie Li, Ziqiang
Cao, and Shuzi Niu. 2017. DailyDialog: A manually
labelled multi-turn dialogue dataset. In Proceedings
of The 8th International Joint Conference on Natural
Language Processing (IJCNLP 2017).
Pierre Lison and Jörg Tiedemann. 2016. OpenSubtitles2016: Extracting large parallel corpora from
movie and TV subtitles. In Proceedings of the Tenth
International Conference on Language Resources
and Evaluation (LREC’16), pages 923–929, Portorož,
Slovenia. European Language Resources Association
(ELRA).
Yinhan Liu, Myle Ott, Naman Goyal, Jingfei Du, Mandar Joshi, Danqi Chen, Omer Levy, Mike Lewis,
Luke Zettlemoyer, and Veselin Stoyanov. 2019.
Roberta: A robustly optimized bert pretraining approach. arXiv preprint arXiv:1907.11692.
Edward Ma. 2019. NLP augmentation.
https://github.com/makcedward/nlpaug.
Tomáš Mikolov, Armand Joulin, Sumit Chopra, Michael
Mathieu, and Marc’Aurelio Ranzato. 2014. Learning
longer memory in recurrent neural networks. arXiv
preprint arXiv:1412.7753.
Robert C. Moore and William Lewis. 2010. Intelligent selection of language model training data. In
Proceedings of the ACL 2010 Conference Short Papers, ACLShort ’10, pages 220–224, Stroudsburg,
PA, USA. Association for Computational Linguistics.
OpenAI, Josh Achiam, Steven Adler, and Sandhini Agarwal et al. 2024. Gpt-4 technical report. arXiv
preprint arXiv:2303.08774. ArXiv:2303.08774 [cs].
U. Orhan, K. E. Hild, D. Erdogmus, B. Roark, B. Oken,
and M. Fried-Oken. 2012. RSVP keyboard: An EEG
based typing interface. In 2012 IEEE International
Conference on Acoustics, Speech and Signal Processing (ICASSP), pages 645–648.
Ondrej Polacek, Adam J. Sporka, and Pavel Slavik.
2017. Text input for motor-impaired people. Universal Access in the Information Society, 16(1):51–72.
Alec Radford, Jeffrey Wu, Rewon Child, David Luan,
Dario Amodei, and Ilya Sutskever. 2018. Language
models are unsupervised multitask learners. OpenAI
blog, 1(8):9.
Colin Raffel, Noam Shazeer, Adam Roberts, Katherine Lee, Sharan Narang, Michael Matena, Yanqi
Zhou, Wei Li, and Peter J. Liu. 2020. Exploring the
limits of transfer learning with a unified text-to-text
transformer. Journal of Machine Learning Research,
21(140):1–67.
Rico Sennrich, Barry Haddow, and Alexandra Birch.
2016. Neural machine translation of rare words with
subword units. In Proceedings of the 54th Annual
Meeting of the Association for Computational Linguistics (Volume 1: Long Papers), page 1715–1725,
Berlin, Germany. Association for Computational Linguistics.
Junxiao Shen, Boyin Yang, John J Dudley, and Per Ola
Kristensson. 2022. Kwickchat: A multi-turn dialogue system for aac using context-aware sentence
generation by bag-of-keywords. In Proceedings of
the 27th International Conference on Intelligent User
Interfaces, IUI ’22, page 853–867, New York, NY,
USA. Association for Computing Machinery.15284Zhiyi Song, Stephanie Strassel, Haejoong Lee, Kevin
Walker, Jonathan Wright, Jennifer Garland, Dana
Fore, Brian Gainor, Preston Cabe, Thomas Thomas,
Brendan Callahan, and Ann Sawyer. 2014. Collecting natural SMS and chat conversations in multiple
languages: The BOLT phase 2 corpus. In Proceedings of the Ninth International Conference on Language Resources and Evaluation (LREC’14), pages
1699–1704, Reykjavik, Iceland. European Language
Resources Association (ELRA).
A. Stolcke. 2002. SRILM – an extensible language modeling toolkit. In International Conference on Spoken
Language Processing, pages 901–904, Denver, CO.
Andreas Stolcke. 1998. Entropy-based pruning of backoff language models. In Proceedings of DARPA
Broadcast News Transcription and Understanding
Workshop.
John Todman, Norman Alm, Jeff Higginbotham, and
Portia File. 2008. Whole utterance approaches in aac.
AAC: Augmentative & Alternative Communication,
24(3):235–254.
Stephanie Valencia, Richard Cave, Krystal Kallarackal,
Katie Seaver, Michael Terry, and Shaun K. Kane.
2023. “the less i type, the better”: How ai language
models can enhance or impede communication for
aac users. In Proceedings of the 2023 CHI Conference on Human Factors in Computing Systems, CHI
’23, page 1–14, New York, NY, USA. Association for
Computing Machinery.
Horabail Venkatagiri. 1999. Efficient keyboard layouts
for sequential access in augmentative and alternative communication. Augmentative and Alternative
Communication, 15(2):126–134.
Keith Vertanen. 2013. A collection of conversational
AAC-like communications. In ASSETS ’13: Proceedings of the ACM SIGACCESS Conference on
Computers and Accessibility.
Keith Vertanen. 2017. Towards improving predictive
AAC using crowdsourced dialogues and partner context. In ASSETS ’17: Proceedings of the ACM
SIGACCESS Conference on Computers and Accessibility (poster), pages 347–348.
Keith Vertanen and Per Ola Kristensson. 2011. The
imagination of crowds: Conversational AAC language modeling using crowdsourcing and large data
sources. In Proceedings of the Conference on Empirical Methods in Natural Language Processing,
EMNLP’11, pages 700–711. ACL.
Keith Vertanen and Per Ola Kristensson. 2014. Complementing text entry evaluations with a composition
task. ACM Transactions on Computer-Human Interaction, 21(2):8:1–8:33.
Keith Vertanen and Per Ola Kristensson. 2021. Mining, analyzing, and modeling text written on mobile
devices. Natural Language Engineering, 27:1–33.
David J Ward and David JC MacKay. 2002. Fast
hands-free writing by gaze direction. Nature,
418(6900):838–838.
Jason Wei, Yi Tay, Rishi Bommasani, Colin Raffel,
Barret Zoph, Sebastian Borgeaud, Dani Yogatama,
Maarten Bosma, Denny Zhou, Donald Metzler, Ed H.
Chi, Tatsunori Hashimoto, Oriol Vinyals, Percy
Liang, Jeff Dean, and William Fedus. 2022. Emergent abilities of large language models. Transactions
on Machine Learning Research. Survey Certification.
Linting Xue, Aditya Barua, Noah Constant, Rami AlRfou, Sharan Narang, Mihir Kale, Adam Roberts,
and Colin Raffel. 2022. ByT5: Towards a token-free
future with pre-trained byte-to-byte models. Transactions of the Association for Computational Linguistics, 10:291–306.
Hussein Yusufali, Stefan Goetze, and Roger K. Moore.
2023. Bridging the Communication Rate Gap: Enhancing Text Input for Augmentative and Alternative Communication (AAC). In HCI International
2023 – Late Breaking Papers, pages 434–452, Cham.
Springer Nature Switzerland.
Hussein Yusufali, Roger K. Moore, and Stefan Goetze.
2024. Refining Text Input For Augmentative and Alternative Communication (AAC) Devices: Analysing
Language Model Layers For Optimisation. In
ICASSP 2024 - 2024 IEEE International Conference on Acoustics, Speech and Signal Processing
(ICASSP), pages 12016–12020. ISSN: 2379-190X.
Susan Zhang, Stephen Roller, Naman Goyal, Mikel
Artetxe, Moya Chen, Shuohui Chen, Christopher Dewan, Mona Diab, Xian Li, Xi Victoria Lin, et al. 2022.
OPT: Open pre-trained transformer language models.
arXiv preprint arXiv:2205.01068.
A Private test set details
A.1 Written test set
We obtained compositions by crowdsourced workers collected by Vertanen and Kristensson (2014).
With the exception of a table of examples and the
text from the AAC condition in Experiment 1, these
novel compositions were never publicly available.
Workers were asked to imagine they were writing
on a mobile device and that they should invent
a “fictitious but plausible message”. We used the
data from condition COMPOSE in Experiment 1,
both conditions in Experiment 2, and the COM-
POSE condition in Experiment 3. See Vertanen and
Kristensson (2014) for further details.
We reviewed workers’ text using a semiautomated process. We expanded abbreviations,
corrected spelling and grammar errors, and split
compositions into sentences. If we could not determine a worker’s intended text, we dropped the15285Are you going to the party?
Do you want to go out for dinner tonight?
Wanna go to taco bell?
I am stuck in a meeting but will call you when I get out.
is it just texting that’s outlawed while driving or talking too
Be sure to watch the meteor shower tonight!
I lost the book.
sarah never called me back
Table 5: Examples of written style communications.
composition. Table 5 shows some examples from
the test set. We subsequently removed these sentences as well as those appearing in tables in Vertanen and Kristensson (2014) from the test set. We
also removed any sentences containing potentially
identifiable information.
We took only the unique compositions from a
given worker. We only used workers who rated
their English ability as native. We assigned each
unique worker an anonymous number. A total of
227 unique workers were in our dataset. We made
a dev and test set from the even and odd numbered
participants respectively. For the purposes of the
experiments conducted here, we lowercased text
and stripped end-of-sentence punctuation, dashes,
and commas. We dropped sentences containing
other characters (e.g. numbers). Our dev and test
sets contained 1,348 lines (8,645 words) and 1,347
lines (9,044 words) respectively.
A.2 Spoken test set
Similar to past work (Venkatagiri, 1999; Vertanen,
2013), we asked people to each write eight different
questions or statements in response to an everyday
communication situation (see Appendix A.4 for
our exact task wording). We did this on Amazon
Mechanical Turk. Each worker received one of the
following situations:
• Opening – Starting a conversation with someone.
• Weather – Talking about the weather.
• Meals – Talking about eating and meals.
• Travel – Talking about trips and traveling.
• Closing – Ending a conversation with someone.
Our study was reviewed by our institutional review board (IRB) and judged to be exempt. Workers were paid $0.80 USD to complete the study.
First, workers completed a consent form and then
created eight unique sentences. While we asked
workers to rate their English ability, we found
this rating was often not accurate. We reviewed
Opening:
How old are you?
Good morning Andrea.
Sorry, Paul i have lots of work in my office.
Weather:
The roads will be hard to see because of the fog.
How many people were killed during Katrina?
Is it going to rain next weekend?
Meal:
Do you prefer chicken or beef?
What’s your favorite kind of meat?
Isn’t this good?
Travel:
What are you looking forward to seeing?
What do you think the food scene is like in Europe?
What is the best place for a vacation in your country?
Closing:
I’m sorry, but I have to go.
It was a pleasure catching up with you.
I will see you around.
Table 6: Examples of spoken style communications in
our five situations.
each worker’s sentences and eliminated 54 workers
who had poor English or misunderstood the task
(e.g. talking about strategies for ending a conversation rather than what you might actually say). This
left us with 246 unique workers who were each
assigned an anonymous number.
Occasionally, we found workers wrote sentences
that were conversational in nature, but did not
match the target situation. We suspect this was
due to worker inattention. We left these sentence in
our test set since we felt they were representative of
conversational messages an AAC user might want
to write. We semi-automatically corrected obvious mistakes in spelling and grammar, and split
compositions comprised of multiple sentences. We
removed any sentences containing potentially identifiable information.
Table 6 shows examples of sentences written
in each situation. In general, we found workers’
compositions were creative and relevant to the target situation. After removing the sentences in Table 6, we made a dev and test set from the even
and odd numbered participants respectively. For
the purposes of the experiments conducted here,
we lowercased text and stripped end-of-sentence
punctuation, dashes, and commas. We dropped sentences containing other characters (e.g. numbers).
Our dev and test sets contained 1,469 lines (11,261
words) and 1,360 lines (10,821 words) respectively.15286A.3 Genre comparison
We found sentences in the written set were shorter
than in the spoken set (6.6 versus 7.8 words per
sentence). We suspect this was due to the mobile
input scenario workers were asked to imagine in
Vertanen and Kristensson (2014). However, given
the input rate limits faced by many AAC users,
the concise nature of the written set may be wellmatched to our target users.
For the spoken set, workers were free to choose
between writing a statement or a question. We classify sentences as a question, statement, or exclamation based on their end-of-sentence punctuation
(i.e. period, question mark, or exclamation point).
In our spoken set, 60% were questions, 36% were
statements, and 4% were exclamations. In the written set, 25% were questions, 43% were statements,
and 11% were exclamations.
A.4 Crowdsourced task
Figure 1 shows the web interface we used to collect AAC-like communications from workers on
Amazon Mechanical Turk. We changed the text at
the top based on the situation as follows:
• Opening – “Imagine you are starting a conversation with a friend, colleague, family member,
or a new acquaintance. Please invent questions
or statements you might use at the start of the
conversation.”
• Weather – “Imagine you are having a conversation about the weather with a friend, colleague,
family member, or a new acquaintance. The
conversation could be about past weather, current weather, or future weather. Please invent
questions or statements you might use in such a
weather-related conversation.”
• Meals – “Imagine you are having a conversation about meals with a friend, colleague, family
member, or a new acquaintance. The conversation
could be about past meals, current meals, or future meals. Please invent questions or statements
you might use in such a meal-related conversation.”
• Travel – “Imagine you are having a conversation
about traveling with a friend, colleague, family
member, or a new acquaintance. The conversation could be about past travels, current travels,
or future travels. Please invent questions or statements you might use in such a travel-related
conversation.”
Figure 1: Worker instructions for our conversational
text collection task (weather communication situation).
• Closing – “Imagine you are speaking with a
friend, colleague, family member, or a new acquaintance. You have reached the point where
you would like to end the conversation. Please
invent questions or statements you might use use
to signal you would like to end the conversation.”
B Additional results and details
B.1 Sentence classifier hyperparameters
We created a model to classify sentences from
C4 and Subtitle on top of DeBERTaV3 (He et al.,
2021). We used a linear learning rate scheduler. We
optimized the model’s hyperparameters via gridsearch as in Table 10 of He et al. (2021).
We searched over batch sizes of [16, 32, 48, 64],
learning rates of [1.5e-5, 2.0e-5, 2.5e-5, 3.0e-5],
warmup steps of [50, 100, 500, 1000], and dropout
of [0.0, 0.1, 0.15]. Weight decay was fixed at 0.01.
Hyperparameter settings were evaluated based on
held out sentences from the the Forum, BOLT,
Daily, and News datasets. Our optimal hyperparameters were: batch size 48, learning rate 2.5e-05,
warmup steps 100, and dropout 0.0.
B.2 Subword hyperparameter tuning
We arrived at our final subword opt-350m model
by first fine-tuning on C4 (threshold 0.90), then
Subtitle (threshold 0.75), then BOLT, and finally
Daily. At each step, we tuned hyperparameters
using Optuna (Akiba et al., 2019) for the equivalent
of 2 GPU days on an A100 GPU.
We used a warmup stable decay (WSD) learning
rate scheduler (Hu et al., 2024). We optimized
with respect to loss on held out sentences in Daily,
BOLT, and Forum. The held out sentences were15287Adaptation data Fine-tuning Dropout Warmup Stable Learning Weight Batch Epochs Tuner Avg
type steps steps rate decay size trials (ppl)
opt-350m base — — — — — — — — — — 2.30
) C4+Sub+BOLT+Daily full 1 .698 751 .586 1.77e-06 .468 1120 1 88 2.19
) C4 full 1 .142 668 .903 2.07e-06 .830 2240 1 80 2.22
) Sub full 2 .245 546 .597 6.71e-08 .025 7840 1 114 2.17
) BOLT full 3 .178 590 .830 2.99e-07 .070 6720 10 849 2.11
) Daily full 4 .435 900 .175 3.53e-04 .288 5600 1 4252 2.10
) BOLT+Daily full 3 .650 49 .638 1.61e-06 .547 2240 1 379 2.09
) BOLT+Daily full 1 .483 905 .009 1.73e-05 .323 6720 9 267 2.07
) C4 LoRA 1 .111 6 .988 1.59e-04 .389 7840 1 75 2.21
) Sub LoRA 2 .003 442 .587 1.39e-06 .368 6720 1 101 2.18
) BOLT+Daily LoRA 3 .385 18 .040 5.20e-05 .535 7840 1 387 2.12
) BOLT+Daily LoRA 1 .524 98 .983 5.73e-04 .898 7840 14 92 2.09
Table 7: Hyperparameters after the specified number of tuning trials during our adaptation process on the opt-350m
LLM. Also shown is the average perplexity on the written and spoken dev sets. We used a sentence selection
threshold of 0.90 for C4 and 0.75 for Subtitle. ) denotes separate fine-tuning steps. + denotes datasets merged in a
single fine-tuning step.
in a 1:1:2 proportion for Daily, BOLT, and Forum
respectively to balance the influence of the written
and spoken text genres. We used an 8-bit version
of the AdamW optimizer.
Fine-tuning examples were single sentences in
mixed case with punctuation and started with the
model’s default start token of </s>. For C4 and
Subtitle, we tuned hyperparameters using four million sentences chosen at random from the C4 0.90
and Subtitle 0.75 datasets to allow a more thorough
search of the parameter space in the allotted time.
Once we selected hyperparameters, we trained the
model using the entire dataset. For Daily and BOLT
we used all the training data from the datasets for
both hyperparameter tuning and training.
We tuned the following hyperparameters, allowing each to vary within the given ranges:
• Dropout — Dropout rate during fine-tuning, [0.0,
0.75].
• Warmup steps — Number of warmup steps in
the WSD learning rate scheduler, [0, 1000].
• Stable — Proportion of the non-warmup steps in
the WSD scheduler that were at a stable (fixed)
learning rate, [0.0, 1.0].
• Learning rate — Learning rate for the stable
phase of the WSD scheduler, [1e-9, 1e-3].
• Weight decay — Weight decay during training,
[1e-4, 1.0].
• Batch size — Training batch size, 1120, 2240,
3360, 4480, 5600, 6720, 7840, or 8960. The batch
size was chosen to maximize the memory utilization of an A100 GPU with 40 GB of memory and
assuming final model training was done using
distributed data parallel (DDP) on four GPUs.
• Epochs — Training epochs over a dataset. We
fixed the training epochs to 1 for C4 and Subtitle.
We allowed the training epochs to be in [1, 20]
for BOLT and Daily.
Table 7 shows the tuned hyperparameter values
as well as the average dev set perplexity for each
step of our adaptation of opt-350m. Each step of
the fine-tuning curriculum lowered the perplexity
on the dev sets. However, we found simply tuning
in one-step on a combination of BOLT and Daily
yielded the best results, reducing the perplexity of
the original opt-350m model from 2.30 to 2.07 (a
10% relative improvement).
We compared the utility of each of our datasets
in isolation in Table 8. We found BOLT performed
the best, followed by Daily, Subtitle 0.75, C4 0.90,
and finally Forum. To show that selecting sentences using our classification model helped, we
fine-tuned on an equivalent random amount of C4
and Subtitle text. As shown in Table 8 (bottom),
random data performed worse than using the threshold selected data. In the case of C4, random data
degraded predictions compared to the base model.
B.3 Potential subword improvements
During our experiments adapting the opt-350m subword LLM, we explored various things that did not
improve the perplexity on our dev sets.15288Written Spoken Avg
(ppl) (ppl) (ppl)
opt-350m base 2.41 2.19 2.30
) Forum 2.37 2.19 2.28
) C4 0.90 2.34 2.09 2.22
) Subtitle 0.75 2.25 2.13 2.19
) Daily 2.25 2.07 2.16
) BOLT 2.19 2.09 2.14
) C4 random 2.47 2.17 2.32
) Subtitle random 2.29 2.15 2.22
Table 8: Perplexity of opt-350m on the dev sets after
single-step fine-tuning on different datasets. ) denotes
separate fine-tuning steps.
Written Spoken Avg
(ppl) (ppl) (ppl)
opt-350m base 2.41 2.19 2.30
) C4 0.95 2.36 2.12 2.24
) C4 0.90 2.34 2.09 2.22
) C4 0.85 2.36 2.11 2.23
) C4 0.90 ) Subtitle 0.80 2.24 2.09 2.17
) C4 0.90 ) Subtitle 0.75 2.24 2.09 2.17
) C4 0.90 ) Subtitle 0.70 2.25 2.10 2.17
Table 9: Perplexity of opt-350m on the dev sets using
different probability thresholds for selecting sentences
from the C4 and Subtitle datasets. ) denotes separate
fine-tuning steps.
Recall that we used n-gram models trained on
data selected from C4 and Subtitle at different
thresholds. We found a threshold of 0.90 for C4
and a threshold of 0.75 for Subtitle provided the
majority of perplexity gains when training an 12gram language model from scratch.
As shown in Table 9, increasing or decreasing
the threshold by 0.05 on C4 increased perplexity.
Subsequent adaptation of the C4 0.90 fine-tuned
model on Subtitle performed best (or the same)
with a threshold of 0.75 versus 0.70 or 0.80.
With recurrent neural network language models
(RNNLMs), interpolating an RNNLM with an ngram model can perform better than either model
in isolation (Mikolov et al., 2014). We computed
the per-character probabilities for our spoken and
written dev sets using 1) the n-gram optimized
for AAC-like text, 2) the original opt-350m LLM,
and 3) our best fine-tuned opt-350m LLM. Using
SRILM (Stolcke, 2002), we found the optimal mixture weights between the n-gram and each LLM
(the unadapted and fine-tuned opt-350m models).
We found the optimal mixture weights with respect
to the spoken and written dev sets independently.
For the unadapted LLM, the optimal mixture had
a perplexity of 2.28, slightly lower than the 2.30
obtained using only opt-350m. For the spoken and
written dev sets, the n-gram model received a mixture weight of 0.12 and 0.29 respectively. However,
the fine-tuned LLM and n-gram mixture model had
the same perplexity of 2.10 as using only the finetuned LLM. For the spoken and written dev sets,
the n-gram model received a mixture weight of
0.04 and 0.03 respectively. Thus it seems the gains
offered by the n-gram model’s in-domain training
data could be obtained instead by fine-tuning a pretrained LLM on in-domain data.
We tried starting each adaptation example with
a unique three-token sequence generated from the
text “ AAC|” rather than the model’s default start
token of “</s>”. We hoped this would allow the
model to associate this token sequence with the
AAC-like text seen during fine-tuning. This might
encourage the model to generate more AAC-like
text when seeing this sequence again at inference
time. We compared the opt-350m model fine-tuned
in a single step on the combined BOLT and Daily
data. We tuned separate model hyperparameters for
each start token sequence. Using a starting context
of “ AAC|” resulted in an average perplexity on
our dev sets of 2.21. Using the default starting context of “</s>” resulted in a slightly lower average
perplexity of 2.18.
B.4 Byte model adaptation
As an alternative to a subword LLM and our search
algorithm, we tested using the ByGPT5 LLM (Belouadi and Eger, 2023). ByGPT5 uses byte instead
of subword tokenization. We fine-tuned ByGPT5
similarly to the subword opt-350m model (see Appendix B.2). Since fine-tuning took around four
times as long compared to opt-350m, we used only
one million sentences when tuning on C4 and Subtitle. This allowed the tuner to explore a similar
number of parameter settings as opt-350m. For the
byte model, the batch size was chosen from 592,
1184, 1776, 2368, 2960, 3552, 4144, or 4736 to
maximize memory use during model fine-tuning.
Table 10 shows the tuned hyperparameters and
the resulting dev set perplexity at each step of
the adaptation curriculum. Predictions improved
markedly compared to the base model, likely due
to the multilingual nature of the original model. We
found that unlike the subword opt-350m model, the
byte model benefited from first fine-tuning on C4
and then subsequently fine-tuning on Subtitle.15289Adaptation data Fine-tuning Dropout Warmup Stable Learning Weight Batch Epochs Tuner Avg
steps steps rate decay size trials (ppl)
ByGPT5 base — — — — — — — — — 2.85
) C4 0.90 1 .458 760 .096 2.77e-04 .382 592 1 88 2.21
) Subtitle 0.75 2 .433 55 .948 6.10e-08 .112 2368 1 156 2.22
) BOLT+Daily 3 .567 142 .109 6.07e-06 .109 6720 2 105 2.14
) BOLT+Daily 1 .701 423 .421 4.83e-05 .514 5600 10 77 2.27
Table 10: Hyperparameters after the specified number of tuning trials during our adaptation process of the ByGPT5
LLM. Also shown is the average perplexity on the dev sets. ) denotes separate fine-tuning steps. + denotes datasets
merged in a single fine-tuning step.
Training data Training Dropout Warmup Stable Learning Weight Batch Examples per Tuner Avg
steps steps rate decay size sentence trials (ppl)
) C4 0.90 1 .583 238 .453 1.08e-04 .642 1120 1 115 2.64
) Subtitle 0.75 2 .361 541 .009 1.01e-06 .438 6720 1 142 2.59
) BOLT 3 .322 955 .807 2.15e-05 .449 5600 5 164 2.55
) Daily 4 .725 340 .569 3.42e-05 .185 8960 4 1881 2.54
) BOLT+Daily 3 .201 13 .535 2.80e-05 .379 3360 20 76 2.47
) BOLT+Daily 1 .011 431 .492 1.19e-04 .687 6720 47 50 2.60
Table 11: Hyperparameters after the specified number of tuning trials during our adaptation process of the opt-350m
model with a classification layer. Also shown is the average perplexity on the dev sets. ) denotes separate training
steps. + denotes datasets merged in a single training step.
B.5 Classifier model adaptation
We experimented with adding a classification layer
to our domain-adapted version of opt-350m (Zhang
et al., 2022). Table 11 (rightmost column) shows
the perplexity on the dev sets at each step in the
curriculum. Because the base classification model
had already received the domain adaptation, the primary aim of this curriculum was to teach the model
the classification task. We did not observe large
perplexity changes between stages as with the byte
model. However, perplexity did decrease slightly
from step-to-step. Similar to during domain adaptation of opt-350m, we found combining BOLT and
Daily in the final step improved performance. We
found it helped to first train the classification model
on C4 and then Subtitle compared to training on
just a combination of BOLT and Daily.
We fine-tuned the opt-350m (Zhang et al., 2022)
classifier model in the same way as the subword
model (Appendix B.2). Since the training examples
chose a location within each sentence (Section 6.3),
we tuned the number of examples created from
each sentence in place of epochs. For the C4 and
Subtitle sets, this was fixed at 1. For BOLT, Daily,
and BOLT+Daily, this was chosen from between
1 and 50, inclusive. Table 11 shows the selected
hyperparameters for each step in the curriculum.
C Search algorithm pseudocode
This work relies heavily on our search algorithm
(Section 3). To convey the algorithm’s structure
more clearly, we produced pseudocode in addition
to the Python code in our supplementary materials. Algorithm 1 describes the vocab hash built
at startup to make the search efficient. The hash
maps each possible text sequence to valid matching
tokens. Algorithm 2 describes the search algorithm
used for each inference.
Data: model vocabulary, valid symbol set
Result: initializes vocabulary hash
valid_vocab ← ∅
vocab_hash ← ∅
foreach token, index ∈ model_vocab do
if each character in token ∈ symbol_set then
/* This is a valid token given our
symbol set. */
append index to valid_vocab
/* Add this token index to the hash
for every key that is a prefix
of the token text. */
for i ← 0 to len(token) do
key ← word[0 to i + 1]
append index to vocab_hash[key]
end
end
end
Algorithm 1: Initializes a hash of model tokens that begin with a given prefix.15290Data: context, valid symbol set, beam width, max completed
Result: probability distribution over symbol set
target_index ← len(context)
last_space ← last index of space in context
tokens ← tokenize(context[0 to last_space])
current_hypos, next_hypos ← empty heap
char_dict ← empty dictionary
completed ← 0
/* Store hypotheses as triples of log prob, token sequence, and character length. */
initial_hypo ← 0.0, tokens, last_space
push initial_hypo to current_hypos
while current_hypos not empty do
sort current_hypos by descending log prob
model_input ← all token seq from current_hypos
log_probs ← log softmax(model_output)
foreach hypo, hypo_index ∈ current_hypos do
/* Determine what context needs to be matched based on hypothesis length. */
remaining ← context[hypo[length]) to end]
/* Narrow the search to only tokens that match the remaining context. */
if len(remaining) = 0 then
search_vocab ← valid_vocab
else
search_vocab ← vocab_hash[remaining]
for i ← 1 to len(remaining) do
pref ix ← remaining[0 to i]
if pref ix is single token then
append tokenize(pref ix) to search_vocab
end
end
end
foreach token ∈ search_vocab do
if hypo[length]) + len(token) > target_index then
/* This hypo has surpassed the context. Add its log prob to the list of log probs
for the first character following the context. */
char ← token[target_index − hypo[length])]
log_prob ← hypo[log_prob] + log_probs[hypo_index][token]
append log_prob to char_dict[char]
increment completed
else
/* Add the extended hypothesis to the heap for the next round. */
new_hypo ← log_prob, (hypo[tokens] + token), (hypo[length] + len(token))
/* Prune the list of hypotheses based on the beam width. */
push next_hypo to next_hypos
if count(next_hypos) > beam_width then
pop hypo with lowest log prob from next_hypos
end
end
end
if completed >= max_completed then
/* Completed maximum number of hypotheses. */
break from outer while loop
end
end
current_hypos ← next_hypos
next_hypos ← empty heap
end
log_probs ← empty dictionary
foreach char ∈ symbol_set do
if char ∈ char_dict then
log_probs[char] ← logsumexp(char_dict[char])
else
log_probs[char] ← −inf
end
end
return softmax(log_probs)
Algorithm 2: Conducts a search using a subword model to produce a probability distribution for
the next character.15291
